\documentclass[a4paper,11pt]{article}

% ============================================================
%  Le Collé — Séance 1 : Déterminants, courbes paramétrées
%  © 2026 Ahmed Amine Ouqlani
% ============================================================

% ---------- Polices ----------
\usepackage[utf8]{inputenc}
\usepackage[T1]{fontenc}
\usepackage[sc]{mathpazo}          % texte et mathématiques : Palatino
\usepackage[scaled=0.88]{helvet}   % titres et repères : linéale
\linespread{1.07}
\usepackage{amsmath,amssymb,amsthm,mathrsfs,stmaryrd}
\usepackage{microtype}

% ---------- Mise en page ----------
\usepackage[a4paper,top=2.3cm,bottom=2.3cm,left=2.4cm,right=2.4cm,
            headheight=14pt,headsep=0.8cm,footskip=1.1cm]{geometry}
\setlength{\parindent}{0pt}
\setlength{\parskip}{0.45em plus 0.15em minus 0.1em}
\raggedbottom
\emergencystretch=1.5em
\clubpenalty=10000 \widowpenalty=10000 \displaywidowpenalty=10000 \brokenpenalty=10000

% ---------- Couleurs : un bleu de nuit, un ocre, des gris ----------
\usepackage{xcolor}
\definecolor{Ink}{HTML}{1C2533}
\definecolor{Bleu}{HTML}{1F3A5F}
\definecolor{BleuMoyen}{HTML}{4A6A8F}
\definecolor{BleuClair}{HTML}{9DB2C9}
\definecolor{Fond}{HTML}{F2F4F7}
\definecolor{Ocre}{HTML}{B07A26}
\definecolor{Brique}{HTML}{A23B2A}
\definecolor{Gris}{HTML}{5E6672}
\definecolor{Filet}{HTML}{C5CCD6}

% ---------- Packages ----------
\usepackage{enumitem}
\setlist{itemsep=0.25em,topsep=0.35em,parsep=0pt}
\setlist[itemize]{label=\textcolor{BleuMoyen}{\small\textbullet},leftmargin=1.4em}
\setlist[enumerate]{leftmargin=1.8em,label=\textsf{\small\color{Bleu}\arabic*.}}
\usepackage{booktabs,array,multicol}
\usepackage{needspace,float}
\usepackage[font=small,labelfont={sf,color=Bleu},labelsep=period,skip=6pt]{caption}
\usepackage{tikz}
\usetikzlibrary{arrows.meta,calc,decorations.markings}
\usepackage{pgfplots}
\pgfplotsset{compat=1.17}
\usepackage{tkz-tab}
\usepackage[most]{tcolorbox}
\usepackage{titlesec,titletoc}
\usepackage{fancyhdr,lastpage,etoolbox}
\usepackage[colorlinks=true,linkcolor=Bleu,urlcolor=Bleu,citecolor=Bleu,
            pdftitle={Le Collé — Propriétés métriques des courbes : sélection de 5 problèmes},pdfauthor={Ahmed Amine Ouqlani}]{hyperref}

\numberwithin{equation}{section}
\numberwithin{figure}{section}
\numberwithin{table}{section}

% ---------- Titres ----------
% Chapitres (sections) : grand numéro, titre en linéale, filet ocre
\titleformat{\section}[display]
  {\normalfont\sffamily\color{Ink}}
  {{\small\color{Gris}\textls[200]{\MakeUppercase{Chapitre}}}\\[-0.2em]%
   {\fontsize{64}{64}\selectfont\color{BleuClair}\thesection}}
  {0.4em}
  {\fontsize{25}{30}\selectfont\bfseries}
  [\vspace{0.5em}{\color{Ocre}\rule{3.2cm}{1.4pt}}]
\titlespacing*{\section}{0pt}{0pt}{1.6em}
\titleformat{name=\section,numberless}[display]
  {\normalfont\sffamily\color{Ink}}{}{0pt}{\fontsize{25}{30}\selectfont\bfseries}
  [\vspace{0.5em}{\color{Ocre}\rule{3.2cm}{1.4pt}}]
\titlespacing*{name=\section,numberless}{0pt}{0pt}{1.6em}

\titleformat{\subsection}
  {\needspace{5\baselineskip}\normalfont\sffamily\Large\bfseries\color{Ink}}
  {\color{Bleu}\thesubsection}{0.8em}{}
\titlespacing*{\subsection}{0pt}{1.8em}{0.5em}
\titleformat{\subsubsection}
  {\needspace{5\baselineskip}\normalfont\sffamily\large\bfseries\color{Ink}}
  {\color{Bleu}\thesubsubsection}{0.7em}{}
\titlespacing*{\subsubsection}{0pt}{1.6em}{0.4em}
\titleformat{\paragraph}[runin]{\normalfont\sffamily\bfseries\color{Ink}}{}{0pt}{}[.]
\titlespacing*{\paragraph}{0pt}{0.9em}{0.6em}
\setcounter{secnumdepth}{3}
\setcounter{tocdepth}{2}

% ---------- Table des matières ----------
\titlecontents{section}[0pt]{\addvspace{0.7em}\sffamily\bfseries\color{Ink}}
  {\makebox[2.2em][l]{\color{Bleu}\thecontentslabel}}{}{\hfill\contentspage}[\addvspace{0.2em}]
\titlecontents{subsection}[2.2em]{\small\linespread{1}\selectfont}
  {\makebox[2.6em][l]{\color{Gris}\thecontentslabel}}{}{\titlerule*[0.7pc]{.}\contentspage}

% ---------- En-têtes et pieds de page ----------
\pagestyle{fancy}
\fancyhf{}
\renewcommand{\sectionmark}[1]{\markboth{\thesection\quad#1}{}}
\fancyhead[L]{\footnotesize\sffamily\color{Gris}\leftmark}
\fancyhead[R]{\footnotesize\sffamily\scshape\color{Gris}Le Collé}
\fancyfoot[C]{\small\sffamily\color{Ink}\thepage}
\fancyfoot[L]{\scriptsize\sffamily\hyperlink{toc}{\color{Filet}sommaire}}
\renewcommand{\headrulewidth}{0.4pt}
\renewcommand{\headrule}{{\color{Filet}\hrule height\headrulewidth width\headwidth}}
\renewcommand{\footrulewidth}{0pt}
\fancypagestyle{plain}{\fancyhf{}\fancyfoot[C]{\small\sffamily\color{Ink}\thepage}\renewcommand{\headrulewidth}{0pt}}

% ---------- Énoncés ----------
\newtheoremstyle{enonce}{0pt}{0pt}{}{}{\sffamily\bfseries\color{Bleu}}{.}{0.5em}
  {\thmname{#1}\thmnumber{~#2}\thmnote{\ \normalfont\sffamily\color{Ink}(#3)}}
\newtheoremstyle{exemple}{0.9em}{0.9em}{}{}{\sffamily\bfseries\color{BleuMoyen}}{.}{0.5em}
  {\thmname{#1}\thmnumber{~#2}\thmnote{\ \normalfont\sffamily\color{Ink}(#3)}}
\newtheoremstyle{remarque}{0.7em}{0.7em}{}{}{\sffamily\itshape}{.}{0.5em}{\thmname{#1}\thmnote{\ (#3)}}
\theoremstyle{enonce}
\newtheorem{definition}{Définition}[section]
\newtheorem{theoreme}[definition]{Théorème}
\newtheorem{proposition}[definition]{Proposition}
\newtheorem{corollaire}[definition]{Corollaire}
\theoremstyle{exemple}
\newtheorem{exemple}[definition]{Exemple}
\theoremstyle{remarque}
\newtheorem*{remarque}{Remarque}

\tcbset{
  defi/.style={enhanced,breakable,colback=white,colframe=white,boxrule=0pt,arc=0pt,
     borderline west={1.2pt}{0pt}{Bleu},left=11pt,right=4pt,top=3pt,bottom=3pt,
     before skip=1.1em,after skip=1.1em},
  thm/.style={enhanced,breakable,colback=Fond,colframe=Fond,boxrule=0pt,arc=0pt,
     borderline west={2.4pt}{0pt}{Bleu},left=11pt,right=8pt,top=6pt,bottom=6pt,
     before skip=1.2em,after skip=1.2em}}
\tcolorboxenvironment{definition}{defi}
\tcolorboxenvironment{theoreme}{thm}
\tcolorboxenvironment{proposition}{thm}
\tcolorboxenvironment{corollaire}{thm}

\renewcommand{\proofname}{Démonstration}
\renewcommand{\qedsymbol}{\textcolor{BleuClair}{$\blacksquare$}}
\makeatletter
\renewenvironment{proof}[1][\proofname]{\par\pushQED{\qed}\normalfont
  \topsep0.5em\relax\trivlist\item[\hskip\labelsep\sffamily\itshape\color{Gris}#1.]\ignorespaces}
  {\popQED\endtrivlist\@endpefalse}
\makeatother

% ---------- Méthode : filets haut et bas, titre en bandeau ----------
\makeatletter\newcommand{\methlabelfix}{\protected@edef\@currentlabel{\thetcbcounter}}\makeatother
\newtcolorbox[auto counter,number within=section]{methode}[2][]{%
  enhanced,unbreakable,colback=white,colframe=white,boxrule=0pt,arc=0pt,
  borderline north={0.9pt}{0pt}{Bleu},borderline south={0.9pt}{0pt}{Bleu},
  left=4pt,right=4pt,top=8pt,bottom=8pt,before skip=1.4em,after skip=1.4em,
  before upper={\methlabelfix\phantomsection\addcontentsline{met}{metentry}{\protect\makebox[2.6em][l]{\thetcbcounter}#2}%
     {\sffamily\bfseries\small\color{white}\colorbox{Bleu}{\strut\ MÉTHODE~\thetcbcounter\ }}%
     \ \ {\sffamily\bfseries\color{Bleu}#2}\par\smallskip},#1}

% ---------- Repères pédagogiques : un filet de couleur, un mot-clé ----------
\newcommand{\definirbloc}[3]{% nom, couleur, libellé
  \newtcolorbox{#1}[1][]{enhanced,breakable,blanker,
    borderline west={1.4pt}{0pt}{#2},left=10pt,right=2pt,top=1pt,bottom=1pt,
    before skip=1em,after skip=1em,
    before upper={{\sffamily\bfseries\footnotesize\color{#2}\textls[120]{\MakeUppercase{#3}}}\quad},##1}}
\definirbloc{reflexe}{Bleu}{Réflexe}
\definirbloc{astuce}{Ocre}{Astuce}
\definirbloc{attention}{Brique}{Attention}
\definirbloc{enpratique}{BleuMoyen}{En pratique}
\definirbloc{plusloin}{Gris}{Pour aller plus loin}
\definirbloc{strategie}{Bleu}{Stratégie}
\definirbloc{lecture}{BleuMoyen}{Lire le schéma}
% compatibilité avec les parties reprises
\newenvironment{garde}{\begin{attention}}{\end{attention}}
\newenvironment{complementHP}{\begin{plusloin}}{\end{plusloin}}
\newenvironment{interpretation}{\begin{enpratique}}{\end{enpratique}}

% ---------- Exercices, corrigés ----------
\newcounter{exo}[section]
\renewcommand{\theexo}{\thesection.\arabic{exo}}
\newcommand{\niveau}[1]{{\color{Ocre}\footnotesize
  \ifcase#1\or$\bullet\circ\circ\circ\circ$\or$\bullet\bullet\circ\circ\circ$\or$\bullet\bullet\bullet\circ\circ$%
  \or$\bullet\bullet\bullet\bullet\circ$\or$\bullet\bullet\bullet\bullet\bullet$\fi}}
\newcommand{\serie}[1]{\par\needspace{8\baselineskip}\vspace{1.5em}%
  \noindent{\sffamily\bfseries\color{Bleu}#1}\par\nobreak\vspace{-0.4em}%
  \noindent{\color{Filet}\rule{\linewidth}{0.4pt}}\par\nobreak\vspace{0.1em}}
\newenvironment{exo}[3]{%
  \par\needspace{5\baselineskip}\vspace{0.75em}\refstepcounter{exo}\label{e:#1}%
  \def\exolab{#1}%
  \noindent{\sffamily\bfseries\color{Bleu}Exercice~\theexo}\quad\niveau{#2}%
  \hfill{\footnotesize\sffamily\color{Gris}#3}\par\nobreak\vspace{0.15em}}%
  {\par\nopagebreak\noindent\hfill{\scriptsize\sffamily\hyperref[c:\exolab]{\color{Gris}corrigé p.~\pageref{c:\exolab}}}\par}
\newenvironment{corrige}[1]{%
  \par\needspace{6\baselineskip}\vspace{1em}\phantomsection\label{c:#1}%
  \noindent{\sffamily\bfseries\color{Bleu}Exercice~\ref{e:#1}}\quad
  {\scriptsize\sffamily\hyperref[e:#1]{\color{Gris}énoncé p.~\pageref{e:#1}}}\par\nobreak\vspace{0.2em}}%
  {\hfill\textcolor{BleuClair}{$\square$}\par}
\newcommand{\etape}[1]{\par\smallskip\noindent{\sffamily\bfseries\small\color{BleuMoyen}#1.}\ }
\newenvironment{solution}{\par\noindent{\sffamily\bfseries\small\color{BleuMoyen}Solution.}\ }{\hfill\textcolor{BleuClair}{$\square$}\par}
\newcommand{\chercher}{\par\smallskip\noindent{\color{Filet}\rule[0.5ex]{1.2cm}{0.4pt}}\
  {\footnotesize\sffamily\itshape\color{Gris}laisser chercher l'étudiant avant de poursuivre}\par\smallskip}

% ---------- Guide de séance ----------

\newcommand{\marge}[1]{\leavevmode\marginpar{\raggedleft\scriptsize\sffamily\color{Gris}\textls[100]{\MakeUppercase{#1}}}\ignorespaces}
\newcommand{\Etape}[4]{%
  \subsection[Étape #1 --- #3 {\color{Gris}(#2)}]{#3}\label{etape:#1}%
  \noindent{\small\sffamily\color{Gris}{\color{Ocre}\bfseries Étape #1}\quad·\quad #2\quad·\quad #4}\par\medskip}
\newcommand{\question}[1]{\par\smallskip\noindent{\sffamily\bfseries\small\color{BleuMoyen}Question.}\ \textit{#1}\par\smallskip}
\newcommand{\rappel}[1]{\par\smallskip\noindent\phantomsection\addcontentsline{rap}{rapentry}{#1}%
  {\sffamily\bfseries\scriptsize\colorbox{Fond}{\color{Gris}\strut RAPPEL}}\ {\sffamily\small\color{Gris}#1.}\ }
\newcommand{\pourquoi}{\par\noindent{\sffamily\bfseries\small\color{Ocre}Pourquoi ?}\ }

% Ouverture de chapitre : \ouverture{objectifs}{prérequis}{durée}
\newcommand{\ouverture}[3]{%
  \begin{tcolorbox}[enhanced,colback=Fond,colframe=Fond,arc=0pt,boxrule=0pt,left=12pt,right=12pt,top=10pt,bottom=10pt,
     before skip=1.2em,after skip=1.6em]
  \begin{minipage}[t]{0.56\linewidth}
  {\sffamily\bfseries\small\color{Bleu}OBJECTIFS}\par\vspace{0.3em}\small #1
  \end{minipage}\hfill
  \begin{minipage}[t]{0.38\linewidth}
  {\sffamily\bfseries\small\color{Bleu}PRÉREQUIS}\par\vspace{0.3em}\small #2
  \ifx\relax#3\relax\else\par\vspace{0.6em}{\sffamily\bfseries\small\color{Bleu}DURÉE}\par\vspace{0.2em}\small #3\fi
  \end{minipage}
  \end{tcolorbox}}
% Fin de chapitre
\newenvironment{essentiel}{\par\needspace{12\baselineskip}\vspace{1.5em}%
  \begin{tcolorbox}[enhanced,unbreakable,colback=white,colframe=white,boxrule=0pt,arc=0pt,
   borderline north={1.4pt}{0pt}{Ocre},left=2pt,right=2pt,top=8pt,bottom=4pt]
  {\sffamily\bfseries\color{Ocre}L'ESSENTIEL DU CHAPITRE}\par\medskip}{\end{tcolorbox}}

% ---------- Notations ----------
\newcommand{\K}{\mathbb{K}}
\newcommand{\R}{\mathbb{R}}
\newcommand{\C}{\mathbb{C}}
\newcommand{\N}{\mathbb{N}}
\newcommand{\Z}{\mathbb{Z}}
\newcommand{\Mn}[1]{\mathcal{M}_{n}(#1)}
\newcommand{\GLn}[1]{\mathrm{GL}_n(#1)}
\newcommand{\LE}{\mathcal{L}(E)}
\newcommand{\rg}{\operatorname{rg}}
\newcommand{\vect}{\operatorname{Vect}}
\newcommand{\Mat}{\operatorname{Mat}}
\newcommand{\Id}{\operatorname{Id}}
\newcommand{\sh}{\operatorname{sh}}
\newcommand{\ch}{\operatorname{ch}}
\renewcommand{\th}{\operatorname{th}}
\DeclareMathOperator{\argsh}{argsh}
\newcommand{\vv}[1]{\overrightarrow{#1}}
\newcommand{\ve}[1]{\vec{#1}}
\newcommand{\ient}[2]{\llbracket #1,#2\rrbracket}
\newcommand{\dt}[1]{\left|\begin{matrix}#1\end{matrix}\right|}
\newcommand{\dis}{\displaystyle}

% ---------- Figures ----------
\pgfplotsset{trig format plots=rad}
\pgfplotsset{courbe/.style={axis lines=middle,axis equal image,
  axis line style={Gris,-{Stealth[length=5pt]}},
  xlabel={$x$},ylabel={$y$},ticklabel style={font=\scriptsize\sffamily,color=Gris},
  every axis x label/.style={at={(ticklabel* cs:1)},anchor=west,font=\small},
  every axis y label/.style={at={(ticklabel* cs:1)},anchor=south,font=\small},
  samples=300,smooth,width=0.6\linewidth}}
\tikzset{arcetude/.style={Bleu,line width=1.25pt},
         reste/.style={BleuClair,line width=1pt},
         tgte/.style={Ocre,line width=0.9pt},
         asy/.style={Gris,dashed,line width=0.6pt},
         pt/.style={circle,fill=Ink,inner sep=1.2pt}}

% ---------- Navigation ----------
\makeatletter
\newcommand*\l@rapentry{\@dottedtocline{1}{0em}{0em}}
\newcommand*\l@metentry{\@dottedtocline{1}{0em}{0em}}
\newcommand{\listemethodes}{\@starttoc{met}}
\newcommand{\listerappels}{\@starttoc{rap}}
\makeatother
% onglets de chapitre sur le bord droit
\usepackage{eso-pic}
\newif\ifonglets \ongletsfalse
\newcommand{\dessineronglet}{%
  \edef\ongcol{\ifnum\value{section}=3 Ocre\else Bleu\fi}%
  \edef\ongnum{\the\value{section}}%
  \pgfmathsetmacro{\ongy}{3.2+2.3*(\ongnum-1)}%
  \begin{tikzpicture}[remember picture,overlay]
  \fill[\ongcol] ([xshift=-0.75cm,yshift=-\ongy cm]current page.north east)
     rectangle ([yshift=-\ongy cm-1.9cm]current page.north east);
  \node[text=white,font=\sffamily\bfseries\large] at ([xshift=-0.37cm,yshift=-\ongy cm-0.95cm]current page.north east)
     {\ongnum};
  \end{tikzpicture}}
\AddToShipoutPictureBG{\ifonglets\ifnum\value{section}>0 \dessineronglet\fi\fi}
% blocs de navigation et d'enseignement
\definirbloc{commencer}{Gris}{Si vous commencez ici}
\definirbloc{horsprogramme}{Brique}{Hors programme}
\newcommand{\HP}{\,{\sffamily\bfseries\scriptsize\textcolor{Brique}{[HP]}}}
\newcommand{\source}[1]{{\footnotesize\sffamily\color{Gris}\,(#1)}}
\newcommand{\coupdepouce}[1]{\par\vspace{-0.2em}\noindent{\footnotesize\sffamily\color{Gris}%
  \textbf{Coup de pouce} (à donner si l'étudiant bloque) : #1}\par}
\newcommand{\pointarret}[1]{\par\vspace{1em}\noindent{\color{Filet}\rule[0.5ex]{\linewidth}{0.4pt}}\par\nobreak
  \noindent{\footnotesize\sffamily\color{Ocre}\bfseries POINT D'ARRÊT POSSIBLE}\quad{\footnotesize\sffamily\color{Gris}#1}\par\vspace{0.6em}}
\newcommand{\demarche}[1]{\noindent{\sffamily\bfseries\small\color{BleuMoyen}Démarche.}\ \textit{#1}\par\smallskip}

% ---------- Corrigés : idées, astuces, pièges, à retenir ----------
\newcommand{\choix}[1]{\par\noindent{\sffamily\bfseries\small\color{BleuMoyen}Choix de la méthode.}\ #1\par\smallskip}
\newcommand{\gain}[1]{\par\noindent{\sffamily\bfseries\small\color{Ocre}Astuce.}\ #1\par}
\newcommand{\piege}[1]{\par\noindent{\sffamily\bfseries\small\color{Brique}Piège.}\ #1\par}
\newtcolorbox{retenirbox}{enhanced,breakable,colback=Fond,colframe=Fond,arc=0pt,boxrule=0pt,
  borderline west={2pt}{0pt}{Ocre},left=8pt,right=6pt,top=4pt,bottom=4pt,before skip=0.7em,after skip=0.5em}
\newcommand{\retenir}[1]{\begin{retenirbox}{\sffamily\bfseries\footnotesize\color{Ocre}À RETENIR}\quad #1\end{retenirbox}}
\newcommand{\TD}[1]{{\color{Ocre}\bfseries TD\,#1}\ ·\ }
\newcommand{\q}[1]{\par\smallskip\noindent{\sffamily\bfseries\color{Bleu}#1}\ }
\providecommand{\iddots}{\mathinner{\mkern1mu\raise1pt\hbox{.}\mkern2mu\raise4pt\hbox{.}\mkern2mu\raise7pt\hbox{.}\mkern1mu}}

\providecommand{\Ker}{\operatorname{Ker}}\renewcommand{\Im}{\operatorname{Im}}\providecommand{\tr}{\operatorname{tr}}\providecommand{\Sp}{\operatorname{Sp}}

% ================================================================ TD : mise en page propre au livret d'exercices
\fancyhf{}
\fancyhead[L]{\footnotesize\sffamily\color{Gris}Sélection · Propriétés métriques des courbes}
\fancyhead[R]{\footnotesize\sffamily\scshape\color{Gris}Le Collé}
\fancyfoot[C]{\small\sffamily\color{Ink}\thepage\,/\,\pageref{LastPage}}
\fancypagestyle{plain}{\fancyhf{}\fancyfoot[C]{\small\sffamily\color{Ink}\thepage\,/\,\pageref{LastPage}}\renewcommand{\headrulewidth}{0pt}}
\makeatletter
\newcommand{\partie}[3]{% numéro, titre, description
  \par\needspace{10\baselineskip}\vspace{1.4em}\phantomsection\label{partie:#1}\pdfbookmark[0]{Partie #1}{partie.#1}%
  {\sffamily\small\color{Gris}\textls[200]{PARTIE #1}}\par\nobreak\vspace{0.1em}
  {\sffamily\bfseries\LARGE\color{Ink}#2}\par\nobreak\vspace{0.3em}
  {\color{Ocre}\rule{2.6cm}{1.2pt}}\par\nobreak\vspace{0.2em}
  {\small\color{Gris}#3}\par\nobreak\vspace{0.6em}}
\newcommand{\exercice}[7]{% numéro, difficulté, titre, origine, notions, identifiant, page du corrigé
  \par\needspace{8\baselineskip}\vspace{1.1em}\phantomsection
  \def\@currentlabel{#1}\label{e:#6}\pdfbookmark[1]{Exercice #1}{e.#6}%
  \noindent{\sffamily\bfseries\large\color{Bleu}Exercice~#1}\quad\niveau{#2}\hfill{\footnotesize\sffamily\color{Gris}#4}\par\nobreak
  \vspace{0.1em}\noindent{\sffamily\bfseries\color{Ink}#3}\par\nobreak
  \vspace{0.1em}\noindent{\footnotesize\sffamily\color{Gris}Notions : #5.\hfill\mbox{corrigé p.~#7}\par}\nobreak
  \vspace{0.35em}}
\newcommand{\renvoi}[1]{\par\nopagebreak\vspace{0.2em}\noindent\hfill{\scriptsize\sffamily\color{Gris}#1}\par}
\newcommand{\place}{\par\vspace{1.6cm plus 0.6cm minus 0.6cm}}
\newcommand{\correction}[4]{% numéro, titre, identifiant, renvoi
  \par\needspace{7\baselineskip}\vspace{1.3em}\phantomsection
  \def\@currentlabel{#1}\label{e:#3}\pdfbookmark[1]{Exercice #1}{c.#3}%
  {\color{Filet}\hrule height 0.4pt}\vspace{0.5em}
  \noindent{\sffamily\bfseries\large\color{Bleu}Exercice~#1}\quad{\sffamily\bfseries\color{Ink}#2}\hfill{\scriptsize\sffamily\color{Gris}#4}\par\nobreak
  \vspace{0.4em}}
\newcommand{\sommaireligne}[3]{\noindent\hbox to 2.2em{\hfill\sffamily\color{Bleu}#1.}\enspace\begin{minipage}[t]{0.8\linewidth}\raggedright#2\end{minipage}\hfill\pageref{e:#3}\par}
\makeatother
\setlength{\parskip}{0.5em plus 0.15em minus 0.1em}

\makeatletter
\newcommand{\probleme}[7]{% numéro, étape, titre, difficulté, origine, durée, identifiant
  \clearpage\phantomsection\label{p:#7}\pdfbookmark[1]{Problème #1 — #2}{p.#7}%
  {\sffamily\small\color{Gris}\textls[150]{PROBLÈME #1\ \ ·\ \ \MakeUppercase{#2}}}\par\nobreak\vspace{0.15em}
  {\sffamily\bfseries\LARGE\color{Ink}#3}\par\nobreak\vspace{0.3em}
  {\color{Ocre}\rule{2.6cm}{1.2pt}}\par\nobreak\vspace{0.2em}
  \noindent{\footnotesize\sffamily\color{Gris}Difficulté~\niveau{#4}\quad·\quad durée indicative : #6\quad·\quad #5}\par\nobreak\vspace{0.3em}}
\newcommand{\fichepb}[3]{% notions, objectifs, prérequis
  {\small\setlength{\parskip}{0.15em}%
  \noindent\hangindent=2.4cm\hangafter=1\makebox[2.4cm][l]{\sffamily\textcolor{Bleu}{Notions}}#1\par
  \noindent\hangindent=2.4cm\hangafter=1\makebox[2.4cm][l]{\sffamily\textcolor{Bleu}{Objectifs}}#2\par
  \noindent\hangindent=2.4cm\hangafter=1\makebox[2.4cm][l]{\sffamily\textcolor{Bleu}{Prérequis}}#3\par}
  \vspace{0.4em}{\color{Filet}\hrule height 0.4pt}\vspace{0.6em}}
\newcommand{\titresection}[2]{% titre, texte
  \clearpage\phantomsection\pdfbookmark[0]{#1}{s.#1}%
  {\sffamily\bfseries\fontsize{24}{28}\selectfont\color{Ink}#1}\par\vspace{0.3em}
  {\color{Ocre}\rule{3.2cm}{1.4pt}}\par\vspace{0.5em}{\small\color{Gris}#2}\par\vspace{0.6em}}
\newcommand{\soustitre}[3]{% numéro, titre, label
  \par\needspace{6\baselineskip}\vspace{1em}\phantomsection\label{#3}%
  {\color{Filet}\hrule height 0.4pt}\vspace{0.5em}
  \noindent{\sffamily\bfseries\large\color{Bleu}Problème~#1}\quad{\sffamily\bfseries\color{Ink}#2}\par\nobreak\vspace{0.3em}}
\makeatother

\begin{document}
\thispagestyle{plain}
{\sffamily\small\color{Gris}\textls[150]{LE COLLÉ · MATHÉMATIQUES PT · CHAPITRE 5}}\par\vspace{0.2em}
{\sffamily\bfseries\fontsize{26}{30}\selectfont\color{Ink}Sélection --- Propriétés métriques}\par\vspace{0.3em}
{\sffamily\Large\color{BleuMoyen}Cinq problèmes à travailler après le cours}\par\vspace{0.4em}
{\color{Ocre}\rule{3.2cm}{1.4pt}}\par\vspace{0.8em}
{\small Ces cinq problèmes sont la sélection principale pour les séances, \emph{après} l'explication du cours.
Ils vont de la mise en application au niveau des devoirs et des concours ; chacun réunit plusieurs résultats du chapitre.
Ils complètent le TD complet et le livre, qu'ils ne remplacent pas.

\textbf{Mode d'emploi.} Les traiter dans l'ordre. Pour chaque problème : lire l'en-tête (notions, objectifs, prérequis),
chercher au brouillon, puis rédiger. En cas de blocage, lire \emph{une seule} indication (page des indications), puis reprendre.
Ne lire la correction qu'une fois le problème terminé ou après un vrai effort : elle explique le \emph{choix de la méthode}
avant les calculs. L'origine de chaque problème est indiquée en tête ; les problèmes « d'après » un concours sont des
adaptations rédigées par Le Collé, pas les énoncés officiels.}
\vspace{1.2em}
{\small\renewcommand{\arraystretch}{1.45}\setlength{\tabcolsep}{4pt}\noindent\begin{tabular}{@{}r>{\raggedright\arraybackslash\sffamily\footnotesize}p{2.9cm}>{\raggedright\arraybackslash}p{6.1cm}ll@{\hspace{6pt}}ccc@{}}
& & & & & {\footnotesize\sffamily énoncé} & {\footnotesize\sffamily indic.} & {\footnotesize\sffamily corrigé}\\\hline
1 & \textcolor{Gris}{Mise en application} & Longueur et courbure d'une boucle & \niveau{3} & 30\,min & \pageref{p:met-boucle} & \pageref{i:met-boucle} & \pageref{c:met-boucle}\\
2 & \textcolor{Gris}{Méthodes classiques} & La cycloïde : longueur, repère de Frenet, développée & \niveau{3} & 45\,min & \pageref{p:sel-m-cycloide} & \pageref{i:sel-m-cycloide} & \pageref{c:sel-m-cycloide}\\
3 & \textcolor{Gris}{Approfondissement} & Développée de l'astroïde & \niveau{4} & 45\,min & \pageref{p:met-astro-dev} & \pageref{i:met-astro-dev} & \pageref{c:met-astro-dev}\\
4 & \textcolor{Gris}{Problème de synthèse} & L'ellipse : courbure aux sommets et développée & \niveau{4} & 1\,h & \pageref{p:sel-m-ellipse} & \pageref{i:sel-m-ellipse} & \pageref{c:sel-m-ellipse}\\
5 & \textcolor{Gris}{Niveau DS / concours} & La deltoïde & \niveau{5} & 1\,h\,15 & \pageref{p:sel-m-2021} & \pageref{i:sel-m-2021} & \pageref{c:sel-m-2021}\\
\end{tabular}}
\probleme{1}{Mise en application}{Longueur et courbure d'une boucle}{3}{Exercice complémentaire · TD complet, exercice~15}{30\,min}{met-boucle}
\fichepb{point double, longueur, courbure, rayon de courbure.}{Appliquer les formules de la longueur et de la courbure, en cherchant d'abord une écriture simple de $\|f'\|$.}{Points doubles ; $\gamma=\frac{\det(f',f'')}{\|f'\|^3}$.}
On considère la courbe $f(t)=\left(t^2,\ t-\dfrac{t^3}3\right)$, $t\in\R$.
\begin{enumerate}
\item Montrer que la courbe a un unique point double, atteint pour $t=\pm\sqrt3$.
\item Montrer que $\|f'(t)\|=1+t^2$ et calculer la longueur de la boucle.
\item Calculer la courbure en $M(t)$, puis le rayon de courbure en $O$.
\end{enumerate}
\par\vspace{0.8em}\noindent\hfill{\scriptsize\sffamily\color{Gris}indications p.~\pageref{i:met-boucle} · correction p.~\pageref{c:met-boucle}}
\probleme{2}{Méthodes classiques}{La cycloïde : longueur, repère de Frenet, développée}{3}{Livret, exercices 1.1 et 1.4, réunis en un problème}{45\,min}{sel-m-cycloide}
\fichepb{angle moitié, longueur, abscisse curviligne, repère de Frenet, courbure, centre de courbure, développée.}{Écrire $f'=v(\cos\alpha,\sin\alpha)$ pour obtenir d'un coup longueur, repère de Frenet, courbure et centre de courbure.}{$1-\cos t=2\sin^2\frac t2$ ; Définitions du repère de Frenet et de la courbure.}
On considère la cycloïde $f(t)=(t-\sin t,\ 1-\cos t)$, $t\in\R$.
\begin{enumerate}
\item Montrer que $f'(t)=2\sin\frac t2\left(\sin\frac t2,\cos\frac t2\right)$. Calculer la longueur d'une arche ($t\in[0,2\pi]$), puis l'abscisse curviligne $s$ d'origine $O=M(0)$ pour $t\in[0,2\pi]$.
\item Pour $t\in\,]0,2\pi[$, déterminer le repère de Frenet $(\ve T,\ve N)$, l'angle $\alpha$ de la tangente et la courbure $\gamma$.
\item Montrer que le centre de courbure est $I(t)=(t+\sin t,\ \cos t-1)$.
\item En déduire que la développée de la cycloïde est une cycloïde, image de la première par une translation.
\end{enumerate}
\par\vspace{0.8em}\noindent\hfill{\scriptsize\sffamily\color{Gris}indications p.~\pageref{i:sel-m-cycloide} · correction p.~\pageref{c:sel-m-cycloide}}
\probleme{3}{Approfondissement}{Développée de l'astroïde}{4}{Exercice complémentaire · TD complet, exercice~23}{45\,min}{met-astro-dev}
\fichepb{repère de Frenet, centre de courbure, développée, similitude.}{Calculer un centre de courbure et reconnaître la développée par un changement de repère.}{Repère de Frenet, courbure ; Formules trigonométriques.}
On considère l'astroïde $f(t)=(\cos^3t,\sin^3t)$, $t\in\left]0,\frac\pi2\right[$.
\begin{enumerate}
\item Déterminer le repère de Frenet et le rayon de courbure en $M(t)$.
\item Montrer que le centre de courbure est $I(t)=\big((2-\cos2t)\cos t,\ (2+\cos2t)\sin t\big)$, ou encore $I(t)=\big(\cos^3t+3\cos t\sin^2t,\ \sin^3t+3\sin t\cos^2t\big)$.
\item On pose $u=x+y$ et $v=y-x$. Montrer que $u=(\cos t+\sin t)^3$ et $v=(\sin t-\cos t)^3$. En déduire que la développée est l'image de l'astroïde par une similitude dont on donnera l'angle et le rapport.
\end{enumerate}
\par\vspace{0.8em}\noindent\hfill{\scriptsize\sffamily\color{Gris}indications p.~\pageref{i:met-astro-dev} · correction p.~\pageref{c:met-astro-dev}}
\probleme{4}{Problème de synthèse}{L'ellipse : courbure aux sommets et développée}{4}{Livret, exercices 1.3 et 2.2, réunis en un problème}{1\,h}{sel-m-ellipse}
\fichepb{angle de la tangente, courbure, normale, enveloppe, développée.}{Justifier la formule de la courbure, puis obtenir la développée comme enveloppe des normales et vérifier la cohérence des deux points de vue.}{Dérivée de $\arctan$, ou de $\tan$ ; Enveloppe d'une famille de droites.}
Soient $a>b>0$ et l'ellipse $f(t)=(a\cos t,b\sin t)$, $t\in\R$. On note $\alpha$ l'angle de la tangente.
\begin{enumerate}
\item Montrer que, en un point régulier, $\alpha'(t)=\dfrac{\det(f',f'')}{\|f'\|^2}$, puis que la courbure vaut $\gamma=\dfrac{\det(f',f'')}{\|f'\|^3}$. Calculer $\gamma(t)$ pour l'ellipse.
\item Calculer le rayon de courbure aux quatre sommets. Où la courbure est-elle maximale ? minimale ?
\item Écrire une équation de la normale en $M(t)$.
\item En déduire la développée de l'ellipse comme enveloppe des normales.
\item Retrouver, à l'aide de la développée, les rayons de courbure aux sommets $(a,0)$ et $(0,b)$.
\end{enumerate}
\par\vspace{0.8em}\noindent\hfill{\scriptsize\sffamily\color{Gris}indications p.~\pageref{i:sel-m-ellipse} · correction p.~\pageref{c:sel-m-ellipse}}
\probleme{5}{Niveau DS / concours}{La deltoïde}{5}{D'après Banque PT 2021, Mathématiques B, partie II (énoncé adapté par Le Collé)}{1\,h\,15}{sel-m-2021}
\fichepb{affixe complexe, symétries, angle moitié, rebroussement, longueur, centre de courbure, développée, similitude.}{Réunir les chapitres 2 et 3 sur une courbe de concours : symétries, point stationnaire, longueur, développée.}{Tout le chapitre 2 ; Longueur, repère de Frenet, centre de courbure.}
On identifie le plan à $\C$ et l'on considère la courbe $\Gamma$ d'affixe $z(t)=2e^{it}+e^{-2it}$, $t\in\R$ ; on note $f(t)=\big(x(t),y(t)\big)$.
\begin{enumerate}
\item \emph{Préliminaires.} Établir les formules $\cos p-\cos q=-2\sin\frac{p+q}2\sin\frac{p-q}2$ et $\sin p+\sin q=2\sin\frac{p+q}2\cos\frac{p-q}2$.
\item \emph{Étude de $\Gamma$.}
\begin{enumerate}
\item Exprimer $x(t)$ et $y(t)$.
\item Montrer que $z\big(t+\frac{2\pi}3\big)=e^{2i\pi/3}z(t)$ et interpréter géométriquement. Montrer que $z(-t)=\overline{z(t)}$. En déduire qu'on peut restreindre l'étude à $\big[0,\frac\pi3\big]$ et expliquer comment obtenir la courbe entière.
\item Montrer que $x'(t)=-4\sin\frac{3t}2\cos\frac t2$ et $y'(t)=4\sin\frac{3t}2\sin\frac t2$. Étudier les variations de $x$ et $y$ sur $\big[0,\frac\pi3\big]$.
\item Donner une équation de la tangente en $M\big(\frac\pi3\big)$ et vérifier qu'elle passe par $A(2,0)$.
\item Préciser la nature du point $M(0)$ et sa tangente.
\item Calculer la longueur totale de $\Gamma$.
\end{enumerate}
\item \emph{Développée.}
\begin{enumerate}
\item Pour $t\in\big]0,\frac{2\pi}3\big[$, déterminer $\ve T$, l'angle $\alpha$, la courbure, le rayon de courbure et $\ve N$, puis montrer que le centre de courbure est $I(t)=(6\cos t-3\cos2t,\ 6\sin t+3\sin2t)$.
\item Montrer que $I\big(t+\frac\pi3\big)$ a pour affixe $3e^{i\pi/3}z(t)$. En déduire que la développée de $\Gamma$ est l'image de $\Gamma$ par une similitude dont on précisera le rapport et l'angle.
\end{enumerate}
\end{enumerate}
\par\vspace{0.8em}\noindent\hfill{\scriptsize\sffamily\color{Gris}indications p.~\pageref{i:sel-m-2021} · correction p.~\pageref{c:sel-m-2021}}
\titresection{Indications}{Une indication par question délicate. N\'en lire qu\'une à la fois, puis reprendre la recherche.}
\soustitre{1}{Longueur et courbure d'une boucle}{i:met-boucle}
\begin{itemize}[leftmargin=1.2em,itemsep=0.25em]\item 1. $x(t)=x(s)$ avec $t\neq s$ impose $s=-t$.
\item 2. Développer $\|f'(t)\|^2$ : c'est un carré parfait.
\item 3. $\gamma=\frac{\det(f',f'')}{\|f'\|^3}$ ; le signe indique le sens dans lequel la courbe tourne.
\end{itemize}
\soustitre{2}{La cycloïde : longueur, repère de Frenet, développée}{i:sel-m-cycloide}
\begin{itemize}[leftmargin=1.2em,itemsep=0.25em]\item 1. $1-\cos t=2\sin^2\frac t2$ et $\sin t=2\sin\frac t2\cos\frac t2$ ; attention au signe de $\sin\frac t2$.
\item 2. $\ve T=(\cos\alpha,\sin\alpha)$ avec $\alpha=\frac\pi2-\frac t2$ ; $\gamma=\frac{d\alpha}{ds}=\frac{\alpha'(t)}{v}$.
\item 3. $I=M+R\ve N$ avec $R=\frac1\gamma$.
\item 4. Poser $t=u+\pi$.
\end{itemize}
\soustitre{3}{Développée de l'astroïde}{i:met-astro-dev}
\begin{itemize}[leftmargin=1.2em,itemsep=0.25em]\item 1. $f'(t)=3\sin t\cos t\,(-\cos t,\sin t)$ : sur $\left]0,\frac\pi2\right[$, le facteur est positif.
\item 2. $I=M+R\ve N$ ; les deux écritures se comparent avec $\cos2t=\cos^2t-\sin^2t$.
\item 3. Développer $(\cos t+\sin t)^3$ ; puis écrire $\cos t+\sin t=\sqrt2\cos\left(t-\frac\pi4\right)$.
\end{itemize}
\soustitre{4}{L'ellipse : courbure aux sommets et développée}{i:sel-m-ellipse}
\begin{itemize}[leftmargin=1.2em,itemsep=0.25em]\item 1. Dériver $\tan\alpha=\frac{y'}{x'}$ ; puis $\gamma=\frac{d\alpha}{ds}$ avec $\frac{ds}{dt}=\|f'\|$.
\item 2. Sommets : $t\in\left\{0,\frac\pi2,\pi,\frac{3\pi}2\right\}$.
\item 3. La normale passe par $M(t)$ et est orthogonale à $f'(t)$.
\item 4. Dériver l'équation de la normale par rapport à $t$ et résoudre le système $2\times2$.
\item 5. Le centre de courbure en un sommet est le point de la développée de même paramètre.
\end{itemize}
\soustitre{5}{La deltoïde}{i:sel-m-2021}
\begin{itemize}[leftmargin=1.2em,itemsep=0.25em]\item 1. Partir de $\cos a\cos b$, $\sin a\sin b$ et $\sin a\cos b$, avec $a=\frac{p+q}2$, $b=\frac{p-q}2$.
\item 2.(b) $e^{-4i\pi/3}=e^{2i\pi/3}e^{-2i\pi}$.
\item 2.(c) Utiliser les formules de la question 1.
\item 2.(e) $f''(0)$ et $f'''(0)$, puis le théorème $(p,q)$.
\item 2.(f) $\|f'\|=4\big|\sin\frac{3t}2\big|$ ; intégrer sur un tiers de la courbe.
\item 3.(a) $f'=4\sin\frac{3t}2\left(-\cos\frac t2,\sin\frac t2\right)$ ; $\alpha=\pi-\frac t2$.
\item 3.(b) Calculer $I\big(t+\frac\pi3\big)$ en affixes et factoriser $3e^{i\pi/3}$.
\end{itemize}
\titresection{Corrections}{Chaque correction commence par le choix de la méthode ; les encadrés résument ce qui resservira. Les renvois « méthode~1.4 du livret » désignent le livret du chapitre.}
\soustitre{1}{Longueur et courbure d'une boucle}{c:met-boucle}
\q{1.} $x(t)=x(s)$ avec $t\neq s$ impose $s=-t$ ; puis $y(t)=y(-t)$ donne $2\left(t-\frac{t^3}3\right)=0$, soit $t=\pm\sqrt3$ ($t\neq0$). Unique point double : $M(\pm\sqrt3)=(3,0)$.

\q{2.} $f'(t)=(2t,\ 1-t^2)$ et $\|f'\|^2=4t^2+1-2t^2+t^4=(1+t^2)^2$, donc $\|f'(t)\|=1+t^2$. La boucle correspond à $t\in[-\sqrt3,\sqrt3]$ :
\[L=\int_{-\sqrt3}^{\sqrt3}(1+t^2)\,dt=2\left(\sqrt3+\frac{3\sqrt3}3\right)=4\sqrt3.\]

\q{3.} $f''=(2,-2t)$, $\det(f',f'')=-4t^2-2(1-t^2)=-2(1+t^2)$, d'où
\[\gamma(t)=\frac{-2(1+t^2)}{(1+t^2)^3}=-\frac2{(1+t^2)^2},\qquad R(t)=-\frac{(1+t^2)^2}2.\]
En $O=M(0)$ : $|R|=\frac12$. (Le signe négatif : la courbe tourne dans le sens indirect.)
\retenir{une longueur se calcule bien quand $\|f'\|^2$ est un carré parfait : c'est ce qu'il faut chercher avant d'intégrer.}
\par\hfill\textcolor{BleuClair}{$\square$}\par
\soustitre{2}{La cycloïde : longueur, repère de Frenet, développée}{c:sel-m-cycloide}
\choix{on factorise $f'$ par l'angle moitié pour l'écrire $f'=v(\cos\alpha,\sin\alpha)$ avec $v>0$ : on lit alors en même temps la vitesse (longueur), l'angle (courbure) et le repère de Frenet.}
\q{1.} $f'(t)=(1-\cos t,\sin t)=\left(2\sin^2\tfrac t2,\ 2\sin\tfrac t2\cos\tfrac t2\right)=2\sin\tfrac t2\left(\sin\tfrac t2,\cos\tfrac t2\right)$, d'où $\|f'(t)\|=2\left|\sin\frac t2\right|=2\sin\frac t2$ sur $[0,2\pi]$. Donc
\[L=\int_0^{2\pi}2\sin\tfrac t2\,dt=\left[-4\cos\tfrac t2\right]_0^{2\pi}=8,\qquad s(t)=\int_0^t2\sin\tfrac u2\,du=4\left(1-\cos\tfrac t2\right).\]
\piege{$\sqrt{4\sin^2\frac t2}=2|\sin\frac t2|$ : sur $[2\pi,4\pi]$ le sinus est négatif ; sans valeur absolue, la deuxième arche aurait une « longueur » $-8$.}

\q{2.} Sur $]0,2\pi[$, $v=2\sin\frac t2>0$, donc $\ve T=\left(\sin\frac t2,\cos\frac t2\right)=(\cos\alpha,\sin\alpha)$ avec $\alpha=\frac\pi2-\frac t2$, et $\ve N=\left(-\cos\frac t2,\sin\frac t2\right)$ (rotation de $\ve T$ d'angle $\frac\pi2$). Alors
\[\gamma=\frac{d\alpha}{ds}=\frac{\alpha'(t)}{v}=\frac{-1/2}{2\sin(t/2)}=-\frac1{4\sin(t/2)}.\]
($\gamma<0$ : parcourue de gauche à droite, la cycloïde tourne vers la droite.)

\q{3.} $R=\frac1\gamma=-4\sin\frac t2$ et
\[I=M+R\ve N=(t-\sin t,1-\cos t)+4\sin\tfrac t2\left(\cos\tfrac t2,-\sin\tfrac t2\right),\]
c'est-à-dire $I=\left(t-\sin t+2\sin t,\ 1-\cos t-4\sin^2\frac t2\right)$,
soit $I=(t+\sin t,\cos t-1)$, car $4\sin^2\frac t2=2(1-\cos t)$.

\q{4.} En posant $t=u+\pi$ : $I=(u-\sin u+\pi,\ -1-\cos u)=(u-\sin u,\ 1-\cos u)+(\pi,-2)$. La développée est la même cycloïde, translatée de $(\pi,-2)$ : ses rebroussements sont au milieu des arches de la première.
\retenir{la factorisation par l'angle moitié donne d'un coup $v$, $\alpha$, $\ve T$, $\ve N$ et $\gamma$. La développée d'une cycloïde est une cycloïde (propriété utilisée par Huygens pour le pendule isochrone).}
\par\hfill\textcolor{BleuClair}{$\square$}\par
\soustitre{3}{Développée de l'astroïde}{c:met-astro-dev}
\q{1.} $f'(t)=3\sin t\cos t\,(-\cos t,\ \sin t)$ avec $3\sin t\cos t>0$ sur $\left]0,\frac\pi2\right[$. Donc $v=3\sin t\cos t$, $\ve T=(-\cos t,\sin t)$, $\ve N=(-\sin t,-\cos t)$. L'angle $\alpha$ de la tangente vérifie $\ve T=(\cos\alpha,\sin\alpha)$ avec $\alpha=\pi-t$, donc $\frac{d\alpha}{dt}=-1$ et
\[\gamma=\frac{d\alpha}{ds}=\frac{-1}{3\sin t\cos t},\qquad R=-3\sin t\cos t.\]
(Le signe négatif : la courbe tourne dans le sens indirect pour ce sens de parcours.)

\q{2.} $I=M+R\ve N=(\cos^3t,\sin^3t)+3\sin t\cos t\,(\sin t,\cos t)=\big(\cos^3t+3\cos t\sin^2t,\ \sin^3t+3\sin t\cos^2t\big)$. Avec $\sin^2t=\frac{1-\cos2t}2$ : $\cos^3t+3\cos t\sin^2t=\cos t\,(\cos^2t+3\sin^2t)=\cos t\,(2-\cos2t)$, et de même $\sin t\,(\sin^2t+3\cos^2t)=\sin t\,(2+\cos2t)$.

\q{3.} Avec $c=\cos t$, $s=\sin t$ : $u=c^3+3cs^2+3c^2s+s^3=(c+s)^3$ et $v=s^3+3sc^2-c^3-3cs^2=(s-c)^3$. Or $c+s=\sqrt2\cos\left(t-\frac\pi4\right)$ et $s-c=\sqrt2\sin\left(t-\frac\pi4\right)$, donc avec $\tau=t-\frac\pi4$ :
\[u=2\sqrt2\cos^3\tau,\qquad v=2\sqrt2\sin^3\tau.\]
Les coordonnées $(u,v)$ sont celles de $I$, au facteur $\sqrt2$ près, dans le repère orthonormé tourné de $\frac\pi4$ : $\left(\frac u{\sqrt2},\frac v{\sqrt2}\right)=2(\cos^3\tau,\sin^3\tau)$. La développée est donc l'image de l'astroïde par la similitude de centre $O$, de rapport $2$ et d'angle $\frac\pi4$.
\par\hfill\textcolor{BleuClair}{$\square$}\par
\soustitre{4}{L'ellipse : courbure aux sommets et développée}{c:sel-m-ellipse}
\choix{l'angle $\alpha$ n'a pas d'expression simple : on dérive $\tan\alpha=\frac{y'}{x'}$, ce qui justifie la formule du déterminant. Pour la développée, plutôt que de calculer $R$ et $\ve N$ en tout point, on enveloppe les normales : il ne reste qu'un système linéaire $2\times2$.}
\q{1.} Là où $x'\neq0$, $\tan\alpha=\frac{y'}{x'}$ ; en dérivant, $(1+\tan^2\alpha)\alpha'=\frac{x'y''-y'x''}{x'^2}$ et $1+\tan^2\alpha=\frac{x'^2+y'^2}{x'^2}$, d'où $\alpha'=\frac{\det(f',f'')}{\|f'\|^2}$ (même calcul avec $\cot\alpha$ là où $y'\neq0$). Comme $\frac{ds}{dt}=\|f'\|$, $\gamma=\frac{d\alpha}{ds}=\frac{\det(f',f'')}{\|f'\|^3}$. Ici $f'=(-a\sin t,b\cos t)$, $f''=(-a\cos t,-b\sin t)$ et $\det(f',f'')=ab$ :
\[\gamma(t)=\frac{ab}{(a^2\sin^2t+b^2\cos^2t)^{3/2}}>0.\]

\q{2.} En $t=0$ ou $\pi$ (sommets $(\pm a,0)$) : $R=\frac{(b^2)^{3/2}}{ab}=\frac{b^2}a$. En $t=\pm\frac\pi2$ (sommets $(0,\pm b)$) : $R=\frac{a^2}b$. Le dénominateur de $\gamma$ est minimal quand $\sin t=0$ (car $b^2<a^2$) : la courbure est maximale aux extrémités du grand axe, minimale à celles du petit axe.
\gain{contrôle : pour $a=b$, on trouve $R=a$ (cercle).}

\q{3.} La normale passe par $M(t)$ et est orthogonale à $f'(t)=(-a\sin t,b\cos t)$ : $-a\sin t\,(x-a\cos t)+b\cos t\,(y-b\sin t)=0$, soit
\[ax\sin t-by\cos t=(a^2-b^2)\sin t\cos t.\]

\q{4.} En dérivant par rapport à $t$ : $ax\cos t+by\sin t=(a^2-b^2)\cos2t$. Le système formé par ces deux équations a pour déterminant $ab(\sin^2t+\cos^2t)=ab\neq0$ ; on obtient
\[x=\frac{a^2-b^2}a\cos^3t,\qquad y=-\frac{a^2-b^2}b\sin^3t.\]
La développée est une astroïde étirée (astroïde de Lamé), avec quatre rebroussements sur les axes.

\q{5.} En $t=0$, le centre de courbure, point de la développée, est $\left(\frac{a^2-b^2}a,0\right)=\left(a-\frac{b^2}a,0\right)$, à la distance $\frac{b^2}a$ du sommet $(a,0)$. En $t=\frac\pi2$ : $\left(0,-\frac{a^2-b^2}b\right)=\left(0,b-\frac{a^2}b\right)$, à la distance $\frac{a^2}b$ du sommet $(0,b)$. On retrouve les rayons de la question~2.
\retenir{$\gamma=\frac{\det(f',f'')}{\|f'\|^3}$ ; développée $=$ enveloppe des normales : équation de la normale, dérivée par rapport au paramètre, système $2\times2$. Les rebroussements de la développée correspondent aux sommets, où la courbure est extrémale.}
\par\hfill\textcolor{BleuClair}{$\square$}\par
\soustitre{5}{La deltoïde}{c:sel-m-2021}
\choix{tout passe par la forme complexe : les symétries se lisent sur $z(t)$, l'angle moitié factorise les dérivées, et la développée se compare à $\Gamma$ par une similitude, c'est-à-dire une multiplication par un nombre complexe.}
\q{1.} Avec $\cos a\cos b=\frac12[\cos(a-b)+\cos(a+b)]$, $\sin a\sin b=\frac12[\cos(a-b)-\cos(a+b)]$ et $\sin a\cos b=\frac12[\sin(a+b)+\sin(a-b)]$, on pose $a=\frac{p+q}2$, $b=\frac{p-q}2$ (donc $a+b=p$, $a-b=q$) : $-2\sin a\sin b=\cos p-\cos q$ et $2\sin a\cos b=\sin p+\sin q$.

\q{2.(a)} $z(t)=(2\cos t+\cos2t)+i(2\sin t-\sin2t)$ : $x(t)=2\cos t+\cos2t$, $y(t)=2\sin t-\sin2t$.

\q{2.(b)} $z\big(t+\frac{2\pi}3\big)=2e^{it}e^{2i\pi/3}+e^{-2it}e^{-4i\pi/3}=e^{2i\pi/3}\big(2e^{it}+e^{-2it}e^{-2i\pi}\big)=e^{2i\pi/3}z(t)$ : $M\big(t+\frac{2\pi}3\big)$ est l'image de $M(t)$ par la rotation de centre $O$ et d'angle $\frac{2\pi}3$, qui laisse donc $\Gamma$ invariante. Par ailleurs $z(-t)=\overline{z(t)}$ : symétrie d'axe $(Ox)$.
Comme $z$ est $2\pi$-périodique, la rotation permet de se limiter à un intervalle de longueur $\frac{2\pi}3$, par exemple $\big[-\frac\pi3,\frac\pi3\big]$, puis la symétrie à $\big[0,\frac\pi3\big]$. On complète par la réflexion d'axe $(Ox)$, puis par les rotations d'angles $\frac{2\pi}3$ et $\frac{4\pi}3$.

\q{2.(c)} $x'(t)=-2\sin t-2\sin2t=-2(\sin t+\sin2t)=-4\sin\frac{3t}2\cos\frac t2$ et $y'(t)=2\cos t-2\cos2t=4\sin\frac{3t}2\sin\frac t2$ (question~1). Sur $\big[0,\frac\pi3\big]$, $\frac{3t}2\in\big[0,\frac\pi2\big]$ : $x'\leq0$ et $y'\geq0$ ; $x$ décroît de $3$ à $\frac12$, $y$ croît de $0$ à $\frac{\sqrt3}2$. Seul $t=0$ annule les deux dérivées.

\q{2.(d)} $M\big(\frac\pi3\big)=\big(\frac12,\frac{\sqrt3}2\big)$ et $f'\big(\frac\pi3\big)=\big(-2\sqrt3,2\big)$. La tangente a pour équation $2\big(x-\frac12\big)+2\sqrt3\big(y-\frac{\sqrt3}2\big)=0$, soit $x+\sqrt3\,y=2$ ; elle passe par $A(2,0)$.

\q{2.(e)} $M(0)=(3,0)$ et $f'(0)=\ve0$. $f''(0)=(-6,0)\neq\ve0$, donc $p=2$ ; $f'''(0)=(0,6)$ et $\det\big(f''(0),f'''(0)\big)=-36\neq0$, donc $q=3$ : \emph{rebroussement de première espèce}, de tangente l'axe $(Ox)$.
\gain{un point stationnaire sur un axe de symétrie a pour tangente cet axe ou sa perpendiculaire.}

\q{2.(f)} $\|f'(t)\|^2=16\sin^2\frac{3t}2\big(\cos^2\frac t2+\sin^2\frac t2\big)$, donc $\|f'(t)\|=4\big|\sin\frac{3t}2\big|$. Sur $\big[0,\frac{2\pi}3\big]$, $\sin\frac{3t}2\geq0$ :
\[\int_0^{2\pi/3}4\sin\tfrac{3t}2\,dt=\tfrac83\Big[-\cos\tfrac{3t}2\Big]_0^{2\pi/3}=\tfrac{16}3.\]
Par la symétrie de rotation, la longueur totale est $3\times\frac{16}3=16$.

\q{3.(a)} Pour $t\in\big]0,\frac{2\pi}3\big[$, $f'(t)=4\sin\frac{3t}2\left(-\cos\frac t2,\sin\frac t2\right)$ avec $v=4\sin\frac{3t}2>0$, donc $\ve T=(\cos\alpha,\sin\alpha)$ avec $\alpha=\pi-\frac t2$, et $\ve N=\left(-\sin\frac t2,-\cos\frac t2\right)$. Alors $\gamma=\frac{\alpha'(t)}{v}=-\frac1{8\sin(3t/2)}$ et $R=-8\sin\frac{3t}2$. Le centre de courbure est
\[I=M+R\ve N=M+8\sin\tfrac{3t}2\left(\sin\tfrac t2,\cos\tfrac t2\right)=M+\big(4(\cos t-\cos2t),\ 4(\sin t+\sin2t)\big),\]
d'après la question~1, soit $I(t)=(6\cos t-3\cos2t,\ 6\sin t+3\sin2t)$.

\q{3.(b)} En affixes, $I(t)=6e^{it}-3e^{-2it}$, donc
\[I\big(t+\tfrac\pi3\big)=6e^{it}e^{i\pi/3}-3e^{-2it}e^{-2i\pi/3}=e^{i\pi/3}\big(6e^{it}-3e^{-2it}e^{-i\pi}\big)=3e^{i\pi/3}\big(2e^{it}+e^{-2it}\big)=3e^{i\pi/3}z(t).\]
Quand $t$ décrit $\R$, $t+\frac\pi3$ aussi : la développée est l'image de $\Gamma$ par la similitude de centre $O$, de rapport $3$ et d'angle $\frac\pi3$.
\retenir{courbe définie par une affixe : symétries par $z(t+\tau)=e^{i\varphi}z(t)$ ou $z(-t)=\overline{z(t)}$ ; longueur et courbure par l'angle moitié.}
\par\hfill\textcolor{BleuClair}{$\square$}\par
\end{document}
